\section{后训练与强化学习}

\subsection{RL 的核心要素}

\textbf{谢天宝（工程视角）}列举了 RL 的四大要素：
\begin{enumerate}
  \item \textbf{Environment}——题面（问题集）+ Sandbox（执行环境）。不同场景要求差异巨大，有些甚至需要 replicate 整个 Internet。
  \item \textbf{Foundation Model}——Mid-training 质量、Cold Start SFT、模型本身的 RL 调性都很重要。
  \item \textbf{Training Infra}——RL 任务的 context 显著长于 SFT；inference engine、context management、agent memory 设计都需要来回迭代。
  \item \textbf{RL 算法}——算法设计又反过来要求 environment 和 training infra 的配合（如 traceback 支持、efficiency 优化）。
\end{enumerate}
这四者高度耦合，涉及的工程细节极多，很难在学术界完整研究，也很难一次性掌握 know-how。

\begin{importantbox}{曹玉：RL 是闭环控制，问题定义比算法更关键}
RL 本质上是\textbf{带反馈信号的闭环控制}，信号链路中每个环节都重要。DeepSeek R1 之所以成功，关键在于\textbf{问题选取与定义的清晰度}——这远比具体算法更决定成败。

当前 RL 面临的现实挑战：
\begin{itemize}
  \item 环境出 bug 概率极高
  \item 问题选取难度大
  \item Reward 获取 sparse
  \item 不同任务间的矛盾冲突仍依赖 human-centric 方式解决
\end{itemize}
RL 正在朝\textbf{专业化方向}发展：做 coding 的人必须成为好的 reward model，做垂类（金融、法律、医疗）也必须懂领域。
\end{importantbox}

\subsection{RL 的 Scaling Law 缺失}

\textbf{薛福昭}：RL 最大的问题之一是\textbf{缺乏开源级别的 Scaling Law}。开源模型发布的是 train 好的 7B/14B/70B，不是 Scaling Curve 上完整可预测下一个 tier 的模型。导致：
\begin{itemize}
  \item 训完小模型不知道大模型会发生什么
  \item 迭代依赖最终的大 Run（又慢又贵）
  \item RL 本身 loss 上蹿下跳，environment 不稳定（package 更新、异步问题），结论非常 noisy
\end{itemize}
一旦建立稳定的 Scaling Law + Leaderboard，RL 的迭代效率会大幅提升。

\subsection{RL 的泛化性问题}

\textbf{曹玉}：RL ``训什么有什么，但不代表其他方面会提升''。当前 RL 参数更新量仍处于``小扰动''范围（极低 LoRA rank 就能打平 full-parameter），导致单任务能力难以自然泛化到其他领域。但 RL 相比 SFT（offline imitation learning），在理论上应有更好的泛化性——关键是 \textbf{RL 的 Scaling 还远未达到 promise 的量级}。

\textbf{刘子伟}：不追求强泛化到完全 unseen task，能在已知 task 矩阵内有一定\textbf{内插能力}就已经足够 impressive。类比 2018 年 Meta Learning 那一波也没找到 general 的泛化方案。Task 和 Environment 的 Scaling 可能是泛化性的关键。

\textbf{薛福昭}：Generalization 的定义不 clear。如果按 FLOPs 计（每个 task sample 1024 条 trace），RL 对指定 task 的泛化性\textbf{应该更强}（比 SFT 每条 sequence 一个 task）。关键瓶颈是无法获得 trillion 级别的 environments。

\begin{knowledgebox}{刘子伟：去伪存真}
过去半年 RL 因 reasoning 大火涌现了大量变种 trick。但最近的工作（如 Just RL）发现，很多 trick 并不有效——\textbf{最朴素、最干净的方案把超参调好就能达到最优}。经历 bubble 后，社区开始区分哪些进展是真的。
\end{knowledgebox}

\subsection{本章小结}
RL 已从``锦上添花''变为提升模型能力的核心工具，但面临 environment 复杂度、scaling law 缺失、泛化性不足三大挑战。DeepSeek R1 的成功源于问题定义的清晰而非算法创新。RL 正在走向专业化。

%% ===== Agentic & Tool Use =====
